\documentclass{beamer}
\usetheme{Madrid}
\usecolortheme{seahorse}
\usepackage{amsmath, amsfonts}
\usepackage{graphicx}
\usepackage{tikz}
\usepackage{qcircuit}
\usepackage{caption}

\title[QML with NN]{Quantum Machine Learning with Neural Networks}
\author{M Hjorth-Jensen}
\institute{University of Oslo}
\date{\today}

\begin{document}

% Title Slide
\begin{frame}
  \titlepage
\end{frame}

% Outline
\begin{frame}{Outline}
  \tableofcontents
\end{frame}

% Section 1
\section{Overview of QML}
\begin{frame}{Quantum Machine Learning Landscape}
  \begin{itemize}
    \item Quantum Data and Classical Data: \( |\psi\rangle \) vs vectors \( \mathbf{x} \in \mathbb{R}^n \)
    \item Categories:
    \begin{itemize}
        \item Quantum-enhanced classical ML
        \item Quantum-native ML algorithms
        \item Hybrid quantum-classical ML
    \end{itemize}
    \item Resource-aware quantum learning models
  \end{itemize}
\end{frame}

% Section 2
\section{Quantum Computing Recap}
\begin{frame}{Quantum States and Gates}
  \begin{itemize}
    \item Qubit: \( |\psi\rangle = \alpha|0\rangle + \beta|1\rangle \), \( |\alpha|^2 + |\beta|^2 = 1 \)
    \item Common Gates:
    \[
        H = \frac{1}{\sqrt{2}}\begin{bmatrix}1 & 1\\ 1 & -1\end{bmatrix}, \quad
        X = \begin{bmatrix}0 & 1\\ 1 & 0\end{bmatrix}
    \]
    \item Entanglement via CNOT and multi-qubit systems
  \end{itemize}
\end{frame}

\begin{frame}{Quantum Circuit Example}
\[
\Qcircuit @C=1em @R=.7em {
  \lstick{|0\rangle} & \gate{H} & \ctrl{1} & \qw \\
  \lstick{|0\rangle} & \qw      & \targ   & \qw
}
\]
\begin{itemize}
  \item This circuit creates a Bell state: \( \frac{1}{\sqrt{2}}(|00\rangle + |11\rangle) \)
  \item Essential for encoding correlations
\end{itemize}
\end{frame}

% Section 3
\section{Neural Networks and Representational Capacity}
\begin{frame}{Classical Neural Networks Review}
  \begin{itemize}
    \item Universal approximation theorem
    \item Layer-wise structure: \( y = \sigma(Wx + b) \)
    \item Deep networks and expressivity
    \item Limitations in high-dimensional feature space
  \end{itemize}
\end{frame}

% Section 4
\section{Quantum Neural Networks (QNN)}
\begin{frame}{What is a QNN?}
  \begin{itemize}
    \item Quantum analog of classical NNs using parameterized quantum circuits (PQC)
    \item Data encoding layer + entangling layer + variational layer
    \item Output is measurement of observables
    \item Trainable parameters: gate angles \( \theta \)
  \end{itemize}
\end{frame}

\begin{frame}{Generic QNN Architecture}
\begin{center}
  \includegraphics[width=0.85\linewidth]{qnn_diagram.png} % Replace with your QNN diagram
\end{center}
\captionof{figure}{Example architecture of a QNN: Feature map + variational layers + measurement}
\end{frame}

% Section 5
\section{Variational Quantum Circuits (VQCs)}
\begin{frame}{Variational Quantum Circuits}
  \begin{itemize}
    \item PQCs parameterized by \( \theta \)
    \item Optimization via classical optimizer (e.g., COBYLA, SPSA, Adam)
    \item Objective: minimize loss \( L(\theta) \) based on expectation values
    \item Cost function: \( C(\theta) = \langle \psi(\theta)| \hat{H} |\psi(\theta)\rangle \)
  \end{itemize}
\end{frame}

\begin{frame}{Circuit Example: Ansatz Layer}
\[
\Qcircuit @C=1em @R=.7em {
  \lstick{|x_1\rangle} & \gate{R_Y(x_1)} & \multigate{1}{Entangle} & \gate{R_Y(\theta_1)} & \qw \\
  \lstick{|x_2\rangle} & \gate{R_Y(x_2)} & \ghost{Entangle}       & \gate{R_Y(\theta_2)} & \qw
}
\]
\begin{itemize}
  \item Hybrid encoding: data \( x_i \) and learnable parameters \( \theta_i \)
  \item Can be scaled for deeper architectures
\end{itemize}
\end{frame}

% Section 6
\section{Training QNNs}
\begin{frame}{Training and Optimization}
  \begin{itemize}
    \item Gradient-based: Parameter Shift Rule
    \[
    \frac{\partial \langle O \rangle}{\partial \theta} =
    \frac{\langle O \rangle_{\theta + \pi/2} - \langle O \rangle_{\theta - \pi/2}}{2}
    \]
    \item Cost landscapes can suffer from barren plateaus
    \item Mitigation via tailored ansatz or local cost functions
  \end{itemize}
\end{frame}

% Section 7
\section{Applications}
\begin{frame}{Applications of QML + NN}
  \begin{itemize}
    \item Quantum-enhanced image and speech classification
    \item Anomaly detection in quantum systems
    \item Quantum generative adversarial networks (QGANs)
    \item Quantum autoencoders for data compression
  \end{itemize}
\end{frame}

% Section 8
\section{Open Challenges}
\begin{frame}{Challenges in QNNs}
  \begin{itemize}
    \item Hardware constraints: decoherence, limited qubits
    \item Barren plateau problem in deep PQCs
    \item Lack of general-purpose quantum optimizers
    \item Scaling to larger problem instances
  \end{itemize}
\end{frame}

% Section 9
\section{Conclusion and Future Work}
\begin{frame}{Future Directions}
  \begin{itemize}
    \item Efficient hybrid architectures
    \item Adaptive circuit learning
    \item Near-term QML benchmarks
    \item Fault-tolerant QML pipelines
  \end{itemize}
\end{frame}


\end{document}
